\section{Các nghiên cứu liên quan}

Phần này trình bày các nghiên cứu liên quan cho các bài toán phát hiện, theo
dõi và đếm phương tiện.

\subsection{Phát hiện phương tiện}

Phát hiện phương tiện là một trong bước đầu tiên để nhận được các thông tin
trong luồng giao thông. Phát hiện nhằm xác định vị trí và phân loại phương
tiện.

Các phương pháp thị giác máy tính cổ điển có thể chia thành ba hướng: chênh
lệch khung hình (frame difference) có ưu điểm về sự tối giản trong thuật toán
và tốc độ tính toán nhanh, nhưng thường chỉ nhận diện được một phần đối tượng
chuyển động do hạn chế của việc đối chiếu sự khác biệt giữa các khung hình liên
tiếp \cite{song2019}. Phương pháp luồng quang học (optical flow) khắc phục được
việc thiếu hụt thông tin bằng cách cung cấp trường chuyển động toàn diện trên
toàn bức ảnh, song khối lượng tính toán khổng lồ lại làm suy giảm đáng kể tốc
độ xử lý thời gian thực \cite{dai2019}. Đối với kỹ thuật trừ nền (background
subtraction), đối tượng được trích xuất dựa trên độ lệch giữa ảnh đầu vào và mô
hình bối cảnh; tuy nhiên, thách thức lớn nhất của kỹ thuật này là sự khó khăn
trong việc duy trì một mô hình nền chuẩn xác trước các biến động môi trường như
thay đổi ánh sáng hay bóng đổ \cite{dai2019}.

Trong quá khứ, các phương pháp chủ yếu dựa vào các đặc trưng được thiết kế thủ
công như SIFT, HOG và Haar-like, sau đó đưa các đặc trưng này vào bộ phân loại
để huấn luyện, ví dụ SVM và Adaboost, v.v.~\cite{dai2019}. Tuy nhiên, phương
pháp này có độ ổn định kém và tốc độ tính toán chậm.

Hiện nay, hướng học sâu dùng mạng nơ-ron tích chập (CNN) khắc phục được nhiều
hạn chế đó. Trong số các mô hình học sâu, nhóm hai giai đoạn (R-CNN, SPPNet,
Fast R-CNN, Faster R-CNN) đạt độ chính xác cao nhờ sinh vùng đề xuất và sử dụng
CNN để chiết xuất đặc trưng, nhưng tốc độ xử lý chậm, khó đáp ứng thời gian
thực. Nhóm một giai đoạn (YOLO, SSD) đạt sự cân bằng tốt hơn giữa độ chính xác
và tốc độ, phù hợp với các hệ thống giám sát giao thông. Các phiên bản YOLO mới
hơn (YOLOv8, YOLOv9, YOLOv10, YOLOv11, YOLOv12, YOLOv26) liên tục cải tiến kiến
trúc, tăng cường cơ chế chú ý và tối ưu hóa hậu xử lý, giúp nâng cao khả năng
phát hiện vật thể và giảm độ trễ.

\subsection{Theo dõi phương tiện}

Mục tiêu của quá trình theo dõi phương tiện là xây dựng quỹ đạo chuyển động của
đối tượng xuyên suốt các khung hình. Đối với bài toán theo dõi đa đối tượng
(MOT), chiến lược chủ đạo là theo dõi dựa trên phát hiện (Detection-Based
Tracking, DBT), trong đó quá trình ghép nối hoàn toàn phụ thuộc vào các hộp bao
do mô hình nhận diện cung cấp. Các thuật toán nền tảng như SORT và DeepSORT đã
định hình chiến lược này bằng cách kết hợp dự đoán động học và đối sánh đặc
trưng ngoại hình. Dựa trên nền tảng đó, các phương pháp hiện đại như ByteTrack
hay OC-SORT tiếp tục tối ưu hóa quá trình liên kết bằng cách tận dụng cả những
hộp bao có độ tin cậy thấp và xử lý tốt hơn các tình huống vật thể bị che
khuất, qua đó giảm thiểu đáng kể hiện tượng nhảy ID (ID switch) và duy trì quỹ
đạo ổn định.

\subsection{Đếm phương tiện}

Đếm phương tiện sử dụng quỹ đạo đã theo dõi để thống kê số lượng, hướng và loại
phương tiện. Các nghiên cứu hiện tại chủ yếu được phân loại thành ba hướng tiếp
cận: phương pháp dựa trên hồi quy (regression-based), phân cụm (cluster-based)
và phát hiện (detection-based) \cite{dai2019}. Phương pháp hồi quy tập trung mô
hình hóa các hàm dự đoán thông qua việc khai thác đặc trưng không gian từ các
vùng ảnh.  Mặc dù đạt hiệu quả cao trong các bối cảnh đông đúc (nơi ranh giới
giữa các đối tượng không rõ ràng), phương pháp này lại gặp khó khăn trước hiện
tượng biến dạng phối cảnh và các đối tượng có kích thước lớn, do chỉ có thể
cung cấp vị trí bao quát của đối tượng \cite{shokri2025}. Phương pháp phân cụm
tiến hành theo dõi các đặc trưng ngoại hình và động học của đối tượng xuyên
suốt chuỗi khung hình nhằm thiết lập quỹ đạo, từ đó tổng hợp và phân cụm các
quỹ đạo này để ước lượng số lượng thực thể. Phương pháp dựa trên phát hiện đang
là phương pháp phổ biến trong bài toán đếm nhờ các kết quả tốt của quá trình
phát hiện và theo dõi phương tiện. Tuy nhiên, yêu cầu tính toán sẽ lớn hơn so
với các phương pháp trước và cần chi phí cao (GPU) để triển khai trên các thiết
bị~\cite{shokri2025}.
